\section{Discovery-Driven Planning: Exposing the Fragile Assumptions}

\subsection{The Illusion of the Traditional Pro-Forma}

An established corporate division can forecast annual revenue by extrapolating 10 years of audited operating history. A startup enterprise cannot. When an early-stage founder presents a 5-year financial pro-forma showing hockey-stick cash flows, they are not presenting financial analysis; they are presenting a bundle of untested hopes disguised as spreadsheet arithmetic.

Rita McGrath and Ian MacMillan developed \textit{discovery-driven planning} to invert this dangerous practice [5]. Instead of pretending to forecast revenues in uncharted territory, the disciplined founder systematically reverse-engineers the operating conditions required to achieve economic solvency, isolates the most fragile assumptions, and designs cheap tests before committing major capital.

\subsection{The Reverse Income Statement}

Begin with the economic threshold required for institutional survival. Suppose an educational founder determines that their small enterprise requires an annual operating contribution of ₹24 lakh after variable delivery expenses to support modest overhead and sustain the founder full-time.

Assume the proposed advanced technical course is priced such that each enrolled student yields ₹24,000 in net variable contribution after mentor compensation, software licenses, and payment processing fees.

The required commercial milestone is immediate:
\begin{equation}
\text{Required Paid Enrollments} = \frac{\text{₹}24,00,000}{\text{₹}24,000} = 100 \text{ students annually}
\end{equation}

\textbf{This 100-student threshold is not a market forecast; it is a mathematical requirement for enterprise survival.}

Now trace the operational chain backward: what physical and behavioral conditions must hold true for 100 engineers to pay and complete this curriculum every year?

\begin{tcolorbox}[
  enhanced,
  colback=BayesSurface!60,
  colframe=BayesBorder,
  boxrule=0.5pt,
  arc=3pt,
  left=3pt,right=3pt,top=3pt,bottom=3pt,
  boxsep=1pt
]
\scriptsize
\begin{tabularx}{\linewidth}{>{\bfseries\raggedright\arraybackslash}p{2.4cm} X}
\toprule
\textbf{Operational Node} & \textbf{Underlying Assumption Requiring Audit} \\
\midrule
Qualified Awareness & At least 2,500 working engineers encounter the technical syllabus. \\
\addlinespace[2pt]
Diagnostic Intent & 20\% (500 engineers) complete a demanding diagnostic assignment. \\
\addlinespace[2pt]
Offer Conversion & 20\% of qualified applicants (100 students) pay ₹35,000 gross tuition. \\
\addlinespace[2pt]
Refund Rate & Refund requests remain strictly below 5\% across 4 quarterly cohorts. \\
\addlinespace[2pt]
Mentor Unit Cost & Grading and code review labor costs do not exceed ₹11,000 per student. \\
\addlinespace[2pt]
Pedagogical Standard & 85\% of students achieve verified competence on independent benchmarks. \\
\bottomrule
\end{tabularx}
\end{tcolorbox}

\subsection{Focusing on the Most Fragile Node}

Inspect the table. Which node represents the most lethal uncertainty?

Curriculum design is rarely the immediate failure mechanism for a skilled engineer. The lethal unknown is conversion: \textit{Will 100 working professionals pay ₹35,000 for an unaccredited curriculum from a brand-new provider?} If your best preliminary signal indicates that conversion among cold leads is 1.5\% rather than 20\%, your business model is off by an order of magnitude. Spending 4 months polishing video lectures while ignoring that massive gap is institutional suicide.

\subsection{Learning Milestones Versus Spending Milestones}

Conventional venture management rewards \textit{spending milestones}: renting office space, hiring 4 engineers, or deploying an advertising campaign.

Discovery-driven planning ties capital release strictly to \textit{learning milestones}:
\begin{itemize}[leftmargin=10pt,itemsep=1pt,topsep=1pt]
  \item \textbf{Spending Milestone:} \textit{Hire 2 assistant instructors by July.}
  \item \textbf{Learning Milestone:} \textit{Demonstrate that a non-founder instructor can evaluate 10 complex student projects using our scoring rubric with 90\% inter-rater agreement and within 45 minutes per project.}
\end{itemize}

The spending milestone measures the depletion of your bank account; the learning milestone measures the acquisition of authentic institutional capability.

\textbf{Never let reverse-engineered arithmetic tempt you into wishful thinking. A reverse income statement clarifies what must be true; it does not make those conditions one bit more likely.}
